\chapter{Content}

\section{Instanzen Generation}
\subsection{Ja - Instanzen}
Bei den Ja-Instanzen handelt es sich um Punktmengen bei denen die \(\degrees\) Funktion Grade liefert welche eine Gradbeschränkte Triangulation ermöglichen. 
Hierbei ist nicht bekannt wie viele mögliche Lösungen es gibt nur das es mindestens eine gibt.
Der Generelle Ablauf fängt mit einer zufälligen Punktmenge an.
Hierbei wird nur die Anzahl der Punkte vorgegeben.
Danach wird mit einem klassischen Verfahren eine Triangulation erstellt.
Von dieser werden dann die Grade und die Knoten gespeichert.
\subsubsection{Delaunay mit zufälligen Flips}
Das erste klassische Verfahren zur Grad bestimmung ist die Delaunay Triangulation in verbindung mit Flips.
Die Delaunay Triangulation wurde im Kapitel~\ref{sec:delaunay_triangulation} bereits vorgestellt.
Für die erstellung der Delaunay Triangulation wird die Bibliothek \texttt{scipy} genutzt.
Auf dieser Triangulation werden dann \(k\) mal die in~\cref{def:flip} Definierten Flips durchgeführt.
Für die durchführung eines Flips zufällig eine Kante \(e_0\) ausgewählt.
Für diese Kante werden zunächst alle Dreiecke gesucht die diese Kante enthalten.
Dies geschied in dem für die beiden Knoten \(x_{0.0}\) und \( x_{0.1}\) der Kante \(e_0\) alle ausgehende Kanten vergleichen werden.
Sollte zwei dieser Kanten den gleichen Knoten  \(v\) haben muss das Tripel (\(x_{0.0}, x_{0.1}, v\)) ein Dreieck sein welches die Kante \(e_0\) enthält.
Sollten am Ende mehr als zwei Dreicke gefunden werden, so müssen die beiden Dreicke ausgewählt werden welche leer sind. 
Hierbei kann es nur zwei lerrer Dreicke geben da es in einer Triangulation nicht zu überschneidungen kommen darf.
Von diesen beiden Dreicken (\(x_{0.0}, x_{0.1}, v_0\)) und (\(x_{0.0}, x_{0.1}, v_1\)) wird dann die Kante \(e_0\) durch die Kante bestehend aus \(v_0\) und \(v_1\) ersetzt.
\subsubsection{Random Kanten}
Bei der zweiten klassischen Triangulation werden zufällig Kanten hinzugefügt.
Hierbei werden zunächst alle möglichen Kanten zwischen den Punkten erstellt.
Bei diesem Schritt werden Kanten ausgeschlossen welche einen anderen Punkt schneiden.
Diese Kanten werden dann in einer zufälligen Reinfolge betrachet.
Bei der auswahl der finalen Kanten wird geprüft ob die Kante eine der zuvor ausgewählten Kanten schneidet.
Sollte dies nicht der Fall sein wird die Kante ausgewählt.
Dieses Verfahren ist maximal da wird alle möglichen Kanten betrachten.

\subsection{Nein - Instanzen}
\begin{itemize}
	\item Move Point liefert meistens nein Instanzen
	\item Zwei Sollgrade von zwei Punkten Tauschen
\end{itemize}
\section{Heuristiken}
\begin{itemize}
	\item eine Triangulation erstellen
	\item den sollGrad möglichst richtig
	\item Werden mit folgender Formel bewertet
\end{itemize}

\begin{align}
 \text{diff}_i &= \left| \degrees(i) - \degree(i) \right| \label{eq:difference} \\
e_i &= \left\{
\begin{array}{ll}
\frac{\degrees(i) - \text{diff}_i}{\degrees(i)} & \degrees(i) \geq \text{diff}_i \\
0 & \, \textrm{sonst} \\
\end{array}
\right. \label{eq:evaluation}\\
Q_T &= \sum_{i=0}^{n} \frac{e_i}{n} \label{eq:total_evaluation}
\end{align}

Die Formel liefert eine Evaluation (\(Q_T\)) für einen Triangulation(\(T\)) mit \(n\) Knoten.
Hierbei gibt \(\degrees(i)\) den gewünschten Grad an der stelle i an.
Dazu passend gibt \(\degree(i)\) den aktuellen Grad an.
In der Brechnung~\eqref{eq:difference} wird die Difference zwischen dem gewünschten und aktuellen Grad berechnet.
In~\eqref{eq:evaluation} wird diese Abweichung Prozentual bewertet.
Am Ende wird in~\eqref{eq:total_evaluation} ein Durchschnitt über alle Prozentualen Werte gebildet.

\subsection{Flips}
\begin{itemize}
	\item Diagonale Flips
	\item k mal Wiederholen
	\begin{itemize}
		\item Alle Knoten durchgehen
		\item wenn falscher Degree auswählen
		\item alle ausgehenden Kanten berechnen
		\item Kante Auswählen \cref{algo:choose_edge}
		\item Flipe Kante wie oben
		\item wenn Gradbeschränkte Triangulation erfolgreich abbrechen
	\end{itemize}
\end{itemize}
\begin{algorithm}[H]
	\caption{CHOOSE EDGE(\(v\))}
	\label{algo:choose_edge}
	\begin{algorithmic}[1]
		\STATE \(E \gets \{ e \mid e \text{ ist ausgehende Kante von } v \}\)
		\WHILE{\(E \neq \emptyset\)}
			\STATE \(e \gets \) zufälliges Element aus \(E\)
			\STATE \(E \gets E \setminus \{e\}\)
			\STATE \(u \gets \text{Knoten von } e~\AND~v \neq u\)
			\IF{\(\degree(u) \neq \degrees(u)\)}
				\RETURN e
			\ELSIF{zu 20 \%}
				\RETURN e
			\ENDIF
		\ENDWHILE
		\RETURN e
	\end{algorithmic}
\end{algorithm}
\subsection{OrTools mit Zielfunktion}
\label{sec:ortools_heu}
TODO infos über OrTools rausuchen
\begin{itemize}
	\item Solver von Google
	\item Solver mit intersection Konstriend und Zielfunktion maximiere~\eqref{eq:evaluation}
	\item erste Konstriend Kanten auf richtige Anzahl
\end{itemize}
Zunächst sei \(E\) die Menge aller Kanten und \(E_P\) die Menge aller möglicher Kanten. Also Kanten die keine Punkte schneiden.
Für jede Kante \(e \in E\) wird eine Bool Variable \(x_e\) erstellt.
Diese Variable ist wahr wenn die Kante \(e\) in der Triangulation enthalten ist.
Für den Überschneidungs Konstriend werden folgende Klauseln dem OrTools Modeln hinzugefügt.
\[
\neg x_{e_1} \lor \neg x_{e_2} \hspace{3cm} \forall e_1, e_2 \in E_P: e_2 \in \intersection(e_1) 
\]
\begin{itemize}
	\item Zielfunktion ist maximiere \(Q_T\) aus~\eqref{eq:evaluation}
	\item Teiler rausgenommen da für max nicht relevant
	\item für die Difference, Betrag und min muss eine Integer Hilfsvarible eingeführt werden.
\end{itemize}
\begin{itemize}
	\item Zweiter ansatz weniger Integer Variablen nutzen
	\item Degree - sollDegree und sollDegree - Degree größer als Hilfsvarible
	\item Summe maximieren
\end{itemize}

\section{Solver}
\begin{itemize}
	\item Exakte Ergebnise Liefern die Triangulation und Gradbeschränkung erfüllen
	\item alle Solver desten ob sum(degree) = 2* Kanten für Triangulation
\end{itemize}
\subsection{OrTools}
\begin{itemize}
	\item gleicher Solver jetzt mit Gradbeschränkung
	\item Kanten Variblen und Intersection Konstriend wie in~\cref{sec:ortools_heu}
	\item Zwei Grundansätze 
	\begin{itemize}
		\item Kanten Maximieren und dann nach Lösung abbrechen welche richtige Anzahl Kanten hat
		\item Kanten Anzahl festsetzen und ohne Zielfunktion lösen
	\end{itemize}	
	\item Ansatz für die Gradbeschränkung
	\item für jeden Knoten die summe der x für ausgehende Kanten varibalen gleich dem Sollgrad
\end{itemize}	
\subsection{SAT mit Kanten}
TODO Rausfinden wie SAT degree gemacht wird
\begin{itemize}
	\item Alle möglichen Kanten als Bool Variable wie OrTools
	\item Intersection Konstriend wie OrTools
	\item Degree sat mit eingebauter Funktion(Infos raussuchen)
	\begin{itemize}
		\item Hilfsveriablen ab max(len(all\_vars), höchste var in solver) + 1
	\end{itemize}
	\item das waren Grundlagen jetzt Verbesserungen
	\item Kanten Anzahl fest setzen
	\item Convexe Hülle festsetzen
	\item Intersection mit Algo Von Michael
	\item Satformel aus Paper von Michael, Fekete und Phillip Alle Edges
	\item Subset degree
	\begin{itemize}
		\item alle Knoten durchgehen
		\item alle Kombinationen der größe len(enges) - (degree -1) bilden
		\item als Klauseln hinzufügen
		\item Als Bild erklären
	\end{itemize}
	\item Bei degree nicht genau degree sondern nur mindestens degree
	\item muss reichen weil sonnst ein anderer Knoten nicht genügend Kanten bekommt
	\item exclude Edges
	\begin{itemize}
		\item alle Kanten betrachten welche die Punkt menge teilen
		\item testen ob anzahl kanten * 2 = sum(degree)
	\end{itemize}
\end{itemize}
\section{SAT mit Dreicken}
\begin{itemize}
	\item Alle möglichen Dreicke erstellen
	\item Spätere Verbesserung nur Leere Dreicke
	\item jedes Dreieck als Bool Variable
	\item Intersection Konstriend wie bei Kanten 
	\item Degree normaler Knoten degree = sum(Dreicke)
	\item Degree auf Convexe Hülle degree - 1 = sum (Dreicke)
\end{itemize}
\section{Theoretische Mehrere Lösungen}
TODO Wo soll das hin
\begin{itemize}
	\item 3 und 4 Knoten ist Gradbeschränkte Triangulation immer eindeutig
	\item ab 8 bis unentlich gibt es mehrere Lösungen
	\item instaze zeigen und erklären wie mehr Punkt hinzufügen
\end{itemize}